\section{From Gram matrices to the slope comparison}\label{sec:frames}
\hypertarget{source-frames}{}

At the point selected in Section~\ref{sec:family-window}, the period maps have different column lattices but a common cuspidal target. A comparison of mean slopes requires a lower bound for the dimension of their joint image. The Gram matrices control that rank question at one
complex place. A column lattice is the integral source of
one period operator. The row is the common cuspidal target
of these operators with the prescribed level and types.
The height and divisibility estimates give
filtration shifts for the global Hermitian lattices. This chapter
proves the linear algebra for integrating those shifts and,
using the local, geometric and dimension estimates proved
below, closes the normalized prime-sum bound. At the chosen complex place, let \(A_i\) map a column of dimension \(r_i\) into the common target of dimension \(D\). Their sum is
\[
 A:\bigoplus_i\mathbb C^{r_i}\longrightarrow\mathbb C^D,\qquad
 (v_i)_i\longmapsto\sum_i A_i v_i,\qquad AA^*=\sum_i A_i A_i^* .
\]
Lemma~\ref{lem:frame-rank} estimates its image, and
Lemma~\ref{lem:frame-slope} compares the source filtrations on that target.

\begin{lemma}[Rank from Gram estimates]\label{lem:frame-rank}
Let \(A_i\) be complex matrices with \(D>0\) rows and \(r_i\)
columns in a finite collection. Let \(R=\sum_i r_i>0\),
\(K_i=A_i^*A_i\) and \(Q=\sum_i A_iA_i^*\), where \(A_i^*\)
is the conjugate transpose. Write \(I\) for an identity
matrix of the indicated size and \(\|\cdot\|_{\rm HS}\)
for the Hilbert--Schmidt norm. Suppose
\[
 \sum_i\|K_i-I\|_{\rm HS}^2\le eR,\qquad
 \operatorname{tr}(Q^2)\le(1+\eta)(\operatorname{tr}Q)^2/D,
 \qquad 0\le e\le1/4,\quad \eta\ge0.
\]
For \(\varepsilon\ge0\), if \(L_i\subseteq\C^{r_i}\) are
subspaces of total codimension
at most \(\varepsilon R\), then
\[
 \operatorname{codim}\left(\sum_i A_iL_i\right)
          \le D(\eta+4\varepsilon+4\sqrt e).
\]
\end{lemma}

\begin{proof}
Let \(P_i\) project orthogonally onto \(L_i\) and let \(P\) project
onto the orthogonal complement of \(\sum_i A_iL_i\).
Since \(PA_iP_i=0\) and orthogonal projection is a contraction,
\[
 \operatorname{tr}(PQ)=\sum_i\|PA_i(I-P_i)\|_{\rm HS}^2
 \le \sum_i\operatorname{tr}((I-P_i)K_i)
 \le(\varepsilon+\sqrt e)R.
\]
The last inequality uses Hilbert--Schmidt Cauchy--Schwarz and
\(\sum_i\|I-P_i\|_{\rm HS}^2\le R\). Applying the same estimate
with the identity matrices gives
\(\operatorname{tr}Q\ge (1-\sqrt e)R\ge R/2\).
Put \(t=\operatorname{tr}(PQ)/\operatorname{tr}Q\).
Then \(0\le t\le1\) and \(t\le2(\varepsilon+\sqrt e)\).
Cauchy--Schwarz gives
\[
 (1-t)^2(\operatorname{tr}Q)^2
  =(\operatorname{tr}((I-P)Q))^2
  \le (D-\operatorname{rk}P)\operatorname{tr}(Q^2).
\]
The trace-ratio hypothesis now yields
\(\operatorname{rk}P/D\le 1-(1-t)^2/(1+\eta)\le\eta+2t\).
\end{proof}

\begin{construction}
\begin{tabular}{@{}l@{\hspace{1.2em}}l@{}}
\toprule
\textit{From the columns} & \textit{On the common target} \\
\midrule
\(L_i\subseteq\mathbb C^{r_i}\) & \(\sum_i A_iL_i\subseteq\mathbb C^D\) \\
\(P_i=\operatorname{proj}_{L_i}\) & \(PA_iP_i=0\) \\
\(K_i=A_i^*A_i,\quad Q=\sum_i A_iA_i^*\) & \(\operatorname{tr}(PQ)\le(\varepsilon+\sqrt e)R\) \\
Second moment of \(Q\) & \(\operatorname{tr}P\le D(\eta+4\varepsilon+4\sqrt e)\) \\
\bottomrule
\end{tabular}
\caption{The projection argument in Lemma~\ref{lem:frame-rank}. Here \(P\) projects onto the orthogonal complement of the surviving image.}\label{fig:analytic-rank}
\end{construction}

\begin{lemma}[Slope estimate]\label{lem:frame-slope}
Let \(N_i\) and \(C\) be Hermitian lattices over a number field,
with normalized Harder--Narasimhan slopes. For such a lattice
\(M\), write \(\mu(M)\) for its mean slope and
\(\mu_{\min}(M),\mu_{\max}(M)\) for its extreme slopes.
Let \(N_{i,\ge t}\)
and \(C_{\ge t}\) denote the rational filtration subspaces with
slopes at least \(t\). Suppose maps \(B_i:N_i\to C\)
and adjustments \(h_i\) satisfy
\[
 B_i(N_{i,\ge t+h_i})\subseteq C_{\ge t}\qquad(t\in\R).
\]
Let \(D_C\) be the target dimension and let \(0\le\delta<1\)
satisfy \(D_C\ge(1-\delta)D\). At one complex place suppose
\(B_i\) is given, in suitable coordinates, by \(A_i\)
as in Lemma~\ref{lem:frame-rank} followed by projection to
the target.
The ranks of \(N_i\) equal the respective \(r_i\).

List the slopes of all \(N_i\), with multiplicities, after
subtracting \(h_i\) from the slopes of \(N_i\).
Assume this list has mean at least \(s\) and every entry is at
most \(s+U\) for some \(U\ge0\). If \(H>0\) and
\(H\ge s-\mu_{\min}(C)\), then the mean slope \(\mu(C)\) satisfies
\[
 s-\mu(C)\le
 \frac{4U\log(1+H/U)+(\eta+4\sqrt e)H}{1-\delta}.
\]
For \(U=0\) the logarithm term is understood by continuity.
\end{lemma}

\begin{proof}
Let \(\lambda_a\) be the adjusted source slopes and let
\(q_N(\tau)=\#\{a:\lambda_a<s-\tau\}/R\).
Since \(s+U-\lambda_a\ge0\) and its mean is at most \(U\),
\[
 (U+\tau)q_N(\tau)
 \le R^{-1}\sum_a(s+U-\lambda_a)\le U .
\]
For \(U=0\) every \(\lambda_a=s\), hence \(q_N=0\).

Apply Lemma~\ref{lem:frame-rank} to the subspaces
\(N_{i,\ge s-\tau+h_i}\) at the chosen embedding. Their projected
images lie in \(C_{\ge s-\tau}\). A surjective projection cannot
increase the codimension of a subspace, so the fraction \(q_C(\tau)\)
of target slopes below \(s-\tau\) satisfies
\[
 q_C(\tau)\le \frac{D}{D_C}\{\eta+4\sqrt e+4q_N(\tau)\},
 \qquad q_C(\tau)=0\quad(\tau\ge H).
\]
Count each target slope \(\lambda\) over the interval
\(0\le\tau<s-\lambda\). The layer integral gives
\[
 s-\mu(C)\le D_C^{-1}\sum_\lambda(s-\lambda)_+
 =\int_0^H q_C(\tau)\,d\tau .
\]
Insert the two bounds and use \(D/D_C\le(1-\delta)^{-1}\).
For \(U>0\) the variable term integrates to \(4U\log(1+H/U)\);
for \(U=0\) it vanishes.
\end{proof}

For the application take \(M=N_{i,\ge t+h_i}\). If its image had nonzero projection to \(C/C_{\ge t}\), the minimum-slope inequality for this quotient map would give
\[
  t\le\mu_{\min}(M)-h_i\le
       \mu_{\max}(C/C_{\ge t})<t .
\]
Here the last strict inequality is the defining property of the Harder--Narasimhan step. This proves the inclusion for every \(t\), including when the image is zero.

The slope polygon over the totally complex coefficient field \(K\) is formed from saturated \(\mathcal O_K\)-sublattices. For each successive quotient in this filtration the squared-norm slope and the underlying integral rank use
\[
 \lambda\longmapsto 2\lambda+c_K,\qquad
 \operatorname{rk}_{\mathbb Z}=[K:\mathbb Q]\operatorname{rk}_K .
\]
Here \(\lambda\) is the logarithmic covolume slope in the Euclidean realization of this quotient, with the sign convention for degrees. Each step of field rank \(r\) occupies \([K:\mathbb Q]r\) integral ranks. The common offset cancels in slope differences, and the common rank factor cancels in the fractions in Lemma~\ref{lem:frame-slope}.  \label{cor:frames-1}

The factor multiplying \(H\) is important. The height limit
has the form \(H\le b x^{c_{HH}}\) at scale \(x\).
We give the compression of the full blocks and the resulting
product bound. Work at the place of the normalized period
matrices. The selected sequence embedding is isometric by the class
metric in Section~\ref{subsec:selected-column-metric} and its
restriction in Section~\ref{subsec:column-orbit-length}.
The common row identity of Section~\ref{sec:transport}
has the following matrix consequence. Denote the full analytic row operator by \(P_u\). The inclusion \(\iota_u\) of the ramified column is isometric, so in adjoint coordinates it satisfies
\[
 \iota_u^*\iota_u=\mathrm{id}\quad\text{on the ramified column}.
\] Let \(\beta_u\ne0\) be the
scalar from unfolding the period, let \(U_u\) be the column
isometry, let \(V\) be the row isometry common to all
sequences, and let \(J\) be the conjugate-linear row embedding.

For a column vector \(\phi\) and row vector \(F\), the
period identity is the equality of bilinear forms
\[
   \mathcal B_u(\phi,F)
    =\beta_u \langle J V F, P_u\iota_u U_u\phi\rangle .
\]
The inner product is conjugate-linear in the first
variable. Thus \(JV\) carries an orthonormal basis of the common
row space to an orthonormal family in the ambient row space.
Extend it to one orthonormal basis of that space. If \(B_u\)
is the arithmetic pairing matrix in rational bases and \(Z,W_u\)
are the row and column changes to orthonormal bases, the restricted analytic matrix has top block
\[
  \beta_u^{-1}Z^{\mathsf t}B_uW_u .
\] The transpose comes from the bilinear
row evaluation. The other rows of the full matrix remain in the
Gram estimate. If the infinite place is represented by the
conjugate embedding, conjugate all entries. This preserves the
Gram inequalities and the isometries.

Let \(r_u^f\) be the full column rank at the sequence index \(u\).
The analytic Gram operator at \(u,u'\) is
\((c_u c_{u'})^{-1/2}P_u^*P_{u'}\), with \(c_u\) the
nonnegative unfolding constant. This constant is positive
for each sequence once the diagonal Gram error is
less than the positive rank \(r_u^f\). Indeed \(c_u=0\) would make the unfolded diagonal operator zero and give
\[
 \text{squared error}=r_u^f\qquad(c_u=0).
\] Use \(c_u^{-1/2}P_u\) as
the full normalized operator. The extra nonzero scalar on
the top block for each sequence preserves the image of every
complex subspace, which is the condition needed in the
slope argument.

Write \(A_u^f\)
for the full blocks, including the positive normalization
scalar on each block, and \(A_u=A_u^f V_u\) for the
ramified blocks, where \(V_u^*V_u=1\).
Write \(R_f=\sum_u r_u^f\) and let \(R\) be the sum of the
ramified ranks. The dimension of the
full target is \(D\) and its reference dimension is \(D_*>0\).

\begin{lemma}[Compression to the ramified columns]\label{lem:gram-compression}
Let \(W\) be the finite index set of the sequences. Write
\(\eta_0,e_0,\eta_1\) for the off-diagonal Gram error,
diagonal Gram error and dimension error, respectively.
Assume the bounds
\[
\begin{split}
 \|(A_u^f)^*A_{u'}^f\|_{\mathrm{HS}}^2
       &\le (1+\eta_0)\,r_u^f r_{u'}^f/D_*
                         \quad (u\ne u'),\\
 \|(A_u^f)^*A_u^f-1\|_{\mathrm{HS}}^2&\le e_0 r_u^f,\qquad
 D\le(1+\eta_1)D_* .
\end{split}
\]
For \(0<\sigma\le1/20\) take
\(-1\le\eta_0\le\sigma/2\), \(\eta_1\le\sigma/2\),
\(\sqrt{e_0}\le\sigma/2\), \(D_*/R_f\le\sigma\), and
\(R\ge(1-d)R_f\) with \(0\le d\le1/4\).
Then the ramified blocks satisfy the diagonal and trace-ratio bounds
in Lemma~\ref{lem:frame-rank} with \(e=2e_0\) and
\(\eta=5\sigma+3\sqrt d\).
\end{lemma}

\begin{proof}
Let \(\Gamma\) be the size parameter used for the number of
sequences and let \(C_0>0\) be the constant in the two rank
bounds below. Multiplying the sequence count by the uniform per-sequence lower
bound gives the rank ratio:
\[
 R_f=\sum_{u\in W}r_u^f\ge |W|\frac{D_*}{C_0\Gamma^2}
       \ge\frac{D_*}{\sigma}.
\]
For the diagonal error, isometric compression is a
Hilbert--Schmidt contraction. Thus the first hypothesis of
Lemma~\ref{lem:frame-rank} holds with \(e=2e_0\).

With \(Q_f=\sum_u A_u^f(A_u^f)^*\) and \(Q=\sum_u A_u A_u^*\),
we have \(0\le Q\le Q_f\), hence
\(\operatorname{tr}Q^2\le\operatorname{tr}Q_f^2\) by the
nonnegative trace of \((Q_f-Q)(Q_f+Q)\).
The denominator in the trace ratio satisfies
\[
 \operatorname{tr}Q\ge(1-\sqrt e)R\ge(1-\sigma)(1-d)R_f .
\]
Expand \(\operatorname{tr}Q_f^2\) into squared Gram blocks and
bound the diagonal blocks by the triangle inequality. This gives
\[
 \frac{D\operatorname{tr}Q^2}{(\operatorname{tr}Q)^2}
 \le
 \frac{(1+\sigma/2)^2+\sigma(1+\sigma/2)^3}
 {(1-\sigma)^2(1-d)^2}
 \le 1+5\sigma+3\sqrt d .
\]

In the last inequality the numerator is at most \(1+9\sigma/4\),
which is at most \((1+5\sigma)(1-\sigma)^2\) on the stated
range. Also \((1-d)^{-2}\le1+4d\) and
\(4(1+5\sigma)d\le3\sqrt d\).
This is the frame compression, with \(\eta=5\sigma+3\sqrt d\).
\label{cor:frames-2}

\end{proof}

\subsubsection*{The upper limit in the slope integral}
\label{par:marked-height-limit}

The preceding rank estimate bounds the integrand in the slope comparison. To integrate it, we also need a cutoff independent of the chosen translate.

Let \(R\) be the direct row space at the refined level and write
\(\bar\kappa\) for the mean logarithm of the period scalar.
The target of the period map is \(R^\vee\).

\begin{proposition}[Uniform upper limit for the slope integral]
\label{prop:marked-frame-height}
For the selected family, use the integral models of
Section~\ref{subsec:division-level-pins} and the period shifts
of Section~\ref{sec:word-division}. There is a middle exponent \(c_{HH}\),
chosen as a function of the outer parameters and \(\epsilon\), for which
the mean \(s\) in Lemma~\ref{lem:frame-slope} satisfies
\begin{equation}\label{eq:marked-frame-height}
          s-\mu_{\min}(R^\vee)\le b x^{c_{HH}} .
\end{equation}
The bound holds at every point of the family above one threshold.
\end{proposition}

\begin{proof}
\proofstep{Apply the sharp row comparison before changing level.}
On the integral models of
Section~\ref{subsec:division-level-pins} the uncut numerator
bound of Section~\ref{subsec:division-length} holds for
every lift at the same point. Thus the projection and
functional comparisons in Section~\ref{sec:word-division}
use all their original place-to-tuple assignments. The
sharp primitive coordinate map takes the arithmetic lattice
onto \(R\) and carries its Gram at every infinite place to
the chosen Gram. It therefore preserves maximal slope:
it takes subspaces bijectively to subspaces and preserves
the covolume of the induced lattice on each one.

At each
embedding the period scalar of this comparison is the one
in the base-level formula. Its log mean is consequently
\(\bar\kappa\) in both. If \(C_f\ge1\) covers the base
functional estimate in Proposition~\ref{prop:descended-row-height},
level passage gives
\[
    \mu_{\max}(R)+\bar\kappa
       \le b x^{C_f}+14\log|G|\le b x^{C_f+1}
       \qquad (x\ge56).
\]
Here \(\log|G|\le4b\) and \((x-1)x^{C_f}\ge56\).
In this application the numerator is needed for the
surjectivity and functional bounds before passing
levels. \label{cor:frame-height-limit-1}

\proofstep{Reduce the upper limit to five pointwise terms.}
For each sequence \(u\) write \(M_u\) for its reference lattice,
\(N_u\) for its shifted lattice, and \(h_u\) for the numerical
height, including the period shift. The column rank bound and the ramified rank estimate give
\(\operatorname{rk}N_u>0\). Use \(\mu(N_u)\le\mu_{\max}(N_u)\)
in \eqref{eq:upper-slopes} and add and subtract \(\bar\kappa\).
For every \(u\) the result is
\[
 \begin{aligned}
 \mu(N_u)-h_u+\mu_{\max}(R)
 &\le \mu_{\max}(M_u)+(G_a+F-E_u)+CmB_0\\
 &\quad -(h_u+\bar\kappa)+(\mu_{\max}(R)+\bar\kappa).
 \end{aligned}
\]
Here \(C\) bounds the same error in every word of the selected
family. Each term on the right needs a common upper bound
before taking the rank-weighted mean.

\proofstep{Bound the two slope costs in terms of the rank parameters.}
Put \(P=1+f+k+m+t+N_m+L\) and \(M=\max(C_t,5)+3\),
where \(C_t\) is the outer rank exponent. The size conditions
give \(P\le x^M\) eventually. Indeed the terms other than
\(N_m\) are bounded by powers of \(x\) from the size bounds
on the host, the ranks and \(L=x^\epsilon\). Substitution of the chain rank into the sum gives the bound
\[
 P=1+f+k+m+t+\bigl(k+m(t+1)\bigr)+L\le 8x^{M-1}\le x^M
 \qquad(x\ge9) .
\]
Also \(B_0=fmbL\), \(a,b',S_c\le Pb\),
\(\log(kv)\le3P\), and \(\log T_u\le C_T' mL\)
by the count on the selected sequences.

In the reference maximal
slope estimate of Proposition~\ref{prop:level-zero-column-slopes}
the three upper costs are fixed multiples of
\[
 |T|a\log(kv),\quad
 f(b'+S_c)\big(k\log(kv)+\log T_u\big),\quad fk^2 L .
\]
Substitute \(|T|,f,k,m,L\le P\) in these three costs.
They are bounded, respectively, by
\[
 3P^3b,\qquad (6+2|C_T'|)P^4b,\qquad P^4 .
\]
For the middle term use
\(k\log(kv)+\log T_u\le(3+|C_T'|)P^2\).
The fixed linear-combination coefficients are absorbed at
large \(x\) using \(P\le x^M\), so the first term costs at
most \(b x^{4M+1}\). The third term \(CmB_0\) has the
same bound because \(mB_0=fm^2bL\le P^4b\).
\label{cor:frame-height-limit-2}

\proofstep{Bound the gain and the forced term before subtracting the centered exponent.}
The unweighted gain at \(p\le X''\) is at most
\(4m|T|a/p\) and the subtracted share is nonnegative.
Above \(X''\) there is no gain. Thus the squared-norm weight
gives
\[
 G_a\le 2\sum_{p\le X''}(\log p)\frac{4m|T|a}{p}
          \le16m|T|a\log(1+X'').
\]
The forced term of one lift is at most \(16P^7 b/p\) to that
cutoff. In fact its type and ramified negative contributions
are bounded above by zero.

On an inert orbit the coefficient
per label is bounded by \(2/p\) using
\(p/(p^2-1)\le2/p\). The sum has at most \(2f\) labels after
averaging around the orbit, each with size at most
\((k-1)4P^5 b\). This includes the shorter orbit with
coefficient \(1/(p^2-1)\). The zero cases contribute nothing.
For the forced term, the logarithmic prime weights give
\[
  \sum_{j=1}^{m}\sum_{p\le X''}(\log p)\frac{16P^7 b}{p}
  \ll m P^7 b\log(1+X'')\ll P^9 b .
\]
Here \(\log(1+X'')\le4P\) and \(m\le P\).
The full-support exponent estimate gives \(-E_u\le CmB_0\),
since each lift's top logarithm contribution is nonnegative.
The two bounds give the stated upper bound \(b x^{9M+2}\)
on the second term for sufficiently large \(x\). \label{cor:frame-height-limit-3}

\proofstep{Estimate the period-free height at each infinite place.}
The fourth term is the negative of the period-free height
in Section~\ref{sec:word-division}. Its three parts are
the negative top size term, the negative sum of kernel
terms, and the positive sum of lower hook terms. The first
costs \(O((1+|C_T'|+|C_{\rm sz}|)P^4b)\). Its sequence-dependent part has at most \(P\) terms of size \(2Pb\log T_u\); the size correction uses \(L_a\le4P\). The respective upper bounds are therefore
\[
 2P^2 b\log T_u,\qquad 4P^3 b L_a|C_{\rm sz}|\le16P^4b|C_{\rm sz}|.
\]
For a kernel at one real place write \(Y=11P^5b\).
The entries of the row weights obey \(|\nu_i|\le Y\)
by addition of the top, short and chain offsets.
Here \(\log Y\le16P\) and \(|\log\Omega|\le7P\).

For every \(\nu_i\), including those of modulus below one, the common lower estimate reads
\[
 \varphi(\nu_i)-\nu_i\log\Omega\ge -Y\log Y-Y-1-7PY .
\] There are at most \(P\) summands. At the same place
the positive hook numerator has at most \(P\) log gamma
terms with arguments in \([2,10P^5b]\).
The lower endpoint uses the gap above the last entry
\(-b_{j-1}\), and the upper endpoint uses the same weight
bounds. Each \(\log\Gamma(y)\le y\log y\) on this range.
The hook denominator consists of differences of the decreasing
integral weight shifted by the index difference. They are
at least one and can be dropped in an upper bound.

The two bounds at one place follow explicitly from \(Y=11P^5b\):
\[
 \begin{aligned}
 P(Y\log Y+Y+1+7PY)&=O(P^7b),\\
 P(10P^5b)\log(10P^5b)&=O(P^7b).
 \end{aligned}
\]
Here the first line bounds the negative kernel terms and the second
bounds the positive hook numerator. The kernel
corrections per lift are bounded by
\(|C_{\rm sz}|(f N_m^2 L_a+fm^3 N_m^2 L)\);
the hook correction is bounded by \(|C_{\rm sz}| f R^2
\log(b\,\operatorname{rd}(E)g_0g_1 R)\), with this
logarithm at most \(6P^2\). After summing over at most
\(P\) places and \(P\) lifts the cost is \(O(P^9b)\).
Consequently \(-(h_u+\bar\kappa)\le b x^{9M+1}\)
eventually. \label{cor:frame-height-limit-4}

\proofstep{Choose the exponent, then take the rank-weighted mean.}
Choose the middle height parameter \(c_{HH}\) at least
\[
     1+\max\{4M+1,9M+2,9M+1,C_f+1\}.
\]
The threshold for it uses only the outer parameters and
\(\epsilon\). Indeed the sharp functional exponent is
obtained from the base functional bound at the middle choice
made as a function of the outer parameters. The four
remaining powers just computed involve only \(C_t\).
Constants multiplying these powers, including the constants
of the upper slope and the size terms, enter the eventual
scale threshold. This permits later constants to depend
on the middle choice.

For \(x\ge9\) the sum of five terms
with exponents at most \(c_{HH}-1\) is at most \(b x^{c_{HH}}\).
With \(r_u=\operatorname{rk}N_u\), multiply by \(r_u\) and
divide the summed inequality by \(\sum_u r_u>0\). This gives
\[
     \frac{\sum_u r_u(\mu(N_u)-h_u)}{\sum_u r_u}
       +\mu_{\max}(R)\le b x^{c_{HH}}.
\]
The first term is \(s\), since the list in
Lemma~\ref{lem:frame-slope} includes slopes with multiplicities.
Hermitian duality gives \(\mu_{\min}(R^\vee)=-\mu_{\max}(R)\).
This gives the bound in the proposition, using the sharp
functional bound and the reference column slope at the same
selected point. \label{cor:frame-height-limit-5}
\end{proof}


\begin{lemma}[Trace error at the height cutoff]\label{lem:trace-at-height}
Use the full Gram estimates on the selected row, the ramified rank
loss of Section~\ref{sec:ramified-column}, and the common exponent
\(C_H=c_{HH}\) from Proposition~\ref{prop:marked-frame-height}.
Uniformly in the retained translates and their sequence indices,
for sufficiently large \(x\),
\begin{equation}\label{eq:normalized-trace-product}
 (\eta+4\sqrt e)x^{C_H}\le6,\qquad
 (\eta+4\sqrt e)H\le6b\quad (0\le H\le bx^{C_H}).
\end{equation}
\end{lemma}

\begin{proof}
Use the middle choice of Proposition~\ref{prop:marked-frame-height}
both for the number of sequences and for the upper limit
\(H=bx^{c_{HH}}\) in the slope integral. Put
\(C_H=c_{HH}\) and \(\sigma=x^{-C_H}\).
The full-block estimates on the selected
records use the errors \(\eta_0,\eta_1,\sqrt{e_0},\sqrt d\)
bounded above by \(E=\exp(-x^{\epsilon/8})\). The ramified
rank loss in Section~\ref{sec:ramified-column} yields this bound
on \(\sqrt d\) for large \(x\) because \(T_u\ge b^{m/2}\)
uniformly over the sequence indices. The normalized bounds and the full
dimension are supplied by Section~\ref{sec:selected-row-gram}
on the linked row.

Take \(x\) so large
that \(E\le \sigma/18\) and \(x^{C_H}\ge20\).
By Lemma~\ref{lem:gram-compression} the resulting errors satisfy
\[
 3\sqrt d+4\sqrt e\le11E\le\sigma,\qquad
 \eta+4\sqrt e\le5\sigma+11E\le6\sigma .
\]
Multiplication by \(x^{C_H}\), or by the allowed \(H\), proves
\eqref{eq:normalized-trace-product}.
These thresholds depend on the fixed exponents and size
constants. They are independent of the sequence index, its coefficient
place, and the point among the retained translates at scale \(x\).
The product in \eqref{eq:normalized-trace-product} is the
term in the arithmetic comparison \eqref{eq:fourier-compare}.

\end{proof}

\subsection{Divisibility and quotient slopes}

The frame estimate counts directions. To use the finite divisibilities one must count each column slope with the multiplicity it has in the quotient.

Fix a nonempty set of columns indexed by \(W\). For \(u\in W\),
let \(M_u,N_u\) be the reference and quotient column lattices.
Average over the columns with weights proportional to \(\rk N_u\).
Thus for a list \((z_u)_{u\in W}\) the notation \(\overline z\)
means
\[
 \overline z =
 \frac{\sum_{u\in W}(\rk N_u)z_u}{\sum_{u\in W}\rk N_u}.
\]
This counts each slope of the concatenated list in
Lemma~\ref{lem:frame-slope} with its multiplicity. Write \(\mu_C=\mu(C)\)
for the mean slope of the common target lattice \(C\), and
let \(h_u\) be the shift used in Lemma~\ref{lem:frame-slope}.
Let \(T_u>0\) be the column size, \(F\) a common shift,
\(L_u\) the loss in the quotient, and \(E_u\) the exponent
correction.

Write \(G_a\) for the arithmetic quantity from the local lattice calculation. The error \(L_P\) is one-sided, in the direction
\[
 S_{\rm prime}\le G_a+L_P .
\]
Let \(S\) be the sum of the top weights, \(D_0\) the
divisibility main term specified below, \(E_D\) the
divisibility error, and \(E_1\) the error in the maximum
slope bound. We will use the following three inequalities simultaneously on these lattices.
The divisibility main term \(D_0\) in the first line is a real number.
In Section~\ref{sec:word-division} the same letter denotes the rank
in the determinant computation supplying that line.
\begin{align}
 \mu_C+h_u-F&\le D_0-S\log T_u+E_D,\label{eq:division-slopes}\\
 \mu(N_u)&\ge\mu(M_u)+G_a+F-E_u-L_u,\label{eq:lower-slopes}\\
 \mu_{\max}(N_u)&\le\mu_{\max}(M_u)+G_a+F-E_u+E_1.
 \label{eq:upper-slopes}
\end{align}
Write \(a_t\) for the top weights, so \(S=\sum_t a_t\).

Let \(f\) be the number of places used for the \(m\) lifts, let
\(w_{j\rho}>0\) be the weight at lift \(j\) and place \(\rho\),
and let \(v>0\) be the scale in the divisibility comparison.
Write \(\ell_P\) for the real logarithmic derivative in the
centering term, \(Q_c\ge0\) for the error in the column count,
and \(\lambda\) for the common term in the centered exponent estimate.
Define \(\Omega\), \(D_0\) and the logarithmic lower bound
\(L_0\) by the following formulas.
\[
 \begin{aligned}
 D_0 &= m\sum_t(a_t\log a_t-a_t-a_t\log\Omega),\\
 \log\Omega&=\gamma+\log(\pi/v)-\ell_P,\\
 \log T_u &\ge L_0
     =\sum_{j=1}^m\frac1f\sum_{\rho=1}^f
               \log(w_{j\rho}v/\pi)+m\lambda-Q_c .
 \end{aligned}
\]
Let \(B_0=fmb\log b\), let \(k\) be the column rank parameter,
and let \(c_E,c_{\rm count}\ge0\) be constants from the local exponent and column level estimates.
\begin{lemma}[Exponent and level spread]\label{lem:exponent-level-spread}
On the family of Sections~\ref{subsec:selected-exponent}
and~\ref{subsec:selected-word-level}, for every pair \(u,u'\)
at a fixed point one has
\begin{equation}\label{eq:slope-word-count-exponent}
 \left|E_u-mS\lambda\right|\le c_E mB_0,
 \qquad
 \left|\log T_u-\log T_{u'}\right|\le c_{\rm count}\log k .
\end{equation}
\end{lemma}
\begin{proof}
Let \(\mathcal S\) be the finite support of the point and let
\(e_w,\omega_w\) be the absolute ramification index and
normalized logarithmic weight at a place \(w\in\mathcal S\).
Write \(E_{u,j}(w)\) for the local exponent contribution.
The normalized finite-place contribution of lift \(j\) is
\[
 s_{u,j}=\sum_{w\in\mathcal S}\omega_w e_w E_{u,j}(w),
\]
In the squared lattice metric \(E_u=2\sum_j s_{u,j}\), so summing
the two-sided deviation at each lift gives
\[
 |E_u-mS\lambda|\le\sum_{j=1}^m |2s_{u,j}-S\lambda|
                     \le c_E mB_0 .
\]

In the count comparison, write \(X_u\)
for the continuous level parameter from the volume formula and
\(\Gamma_*\ge1\) for the common size correction. The level choice
lies in the interval from \(X_u\Gamma_*^{-2/(fk)}\) to that
number times \(1+1/(fk)\). Hence
\(\log T_u-\log X_u+2\log\Gamma_* /(fk)\in[0,1/(fk)]\).
The volume estimate therefore gives the second inequality with \(c_{\rm count}\ge c_5+1\). Indeed, on the eventual scale,
\[
 |\log X_u-\log X_{u'}|+1/(fk)\le c_5\log k+\log k .
\] Section~\ref{subsec:selected-word-level} gives this volume estimate
and the centered lower bound with \(\lambda=\log(g_0g_1)\) and
\(Q_c=O(L)\). Section~\ref{subsec:selected-exponent} proves the exponent comparison on the same point by summing both type fibres and the absolute off-type cost.
\label{cor:frames-3}
\end{proof}
Thus \(E_0=c_E mB_0\) bounds the error in the mean. To compare that mean with the prime sum, separate the centering term in the logarithms.
Expanding the logarithms yields
\[
 D_0-SL_0+mS\lambda=S_{\mathrm{cent}}+SQ_c,\qquad
 S_{\mathrm{cent}}=m\sum_t a_t(\log a_t-1-\gamma+\ell_P)
                 -S\sum_{j=1}^m\frac1f\sum_{\rho=1}^f\log w_{j\rho}.
\]
This is the centering of the prime-power sum. In particular the
logarithms of the chain weights have coefficient \(-S\).

The two slope estimates for \(M_u\) apply to the level-zero
lattice of the selected sequence. Write \(A=|T|a\)
and let \(c_L L\) be the logarithmic parameter in the two
column bounds. The lower and upper calculations can have
different constants and thresholds. In the notation of the
weights for the sequences their conclusions are
\begin{align*}
 -\mu(M_u) &\le
  \frac{\log T_u}{k}\sum_{t\in T}a_t(u)
  +c_-\bigl(A\log k+fk(b'+S_c)\log k+
        fS_c\log T_u+fk^2 c_L L\bigr),\\
 \mu_{\max}(M_u)&\le
  c_+\bigl(A\log(kv)+f(b'+S_c)(k\log(kv)+\log T_u)
          +fk^2 c_L L\bigr).
\end{align*}
The lower estimate starts from the determinant slope of
the integral Weyl lattice with its character scale. The
upper estimate uses short dual vectors on the class lattices
and a highest-weight norm bound. The finite and infinite
coordinate identifications and both slope proofs are in
Section~\ref{subsec:selected-column-metric}. The local quotients and the ideal
shift of this \(M_u\) are those of
Section~\ref{subsec:column-orbit-length}.
\label{cor:frames-4}

\begin{lemma}[Uniform reference slopes]\label{lem:reference-column-slopes}
Apply the lower and upper bounds of
Proposition~\ref{prop:level-zero-column-slopes} to the level-zero
lattice \(M_u\) of each selected word. Then the selected size bounds give
\[
 -\mu(M_u)\le C_M A L_a(1+m/k),\qquad
 \mu_{\max}(M_u)\le C_M A L_a
\]
at every column beyond one threshold, where
\(C_M=2C_T+\max(0,c_-,c_+)(12+C_T+c_L)\).
\end{lemma}
\begin{proof}
Use \(a_t(u)\le2a\), \(\log T_u\le C_TmL\),
\(fk(b'+S_c),fk^2\le A\), \(m\le k\), and \(L\le L_a\).
In the common size bounds \(\log k\le5L_a\) and
\(\log(kv)\le6L_a\). The first term in the lower estimate costs at
most \(2C_T A(m/k)L_a\). Each parenthesis in the two estimates costs
at most \((12+C_T+c_L)A L_a\), by the preceding inequalities.
The claimed bounds follow. Once the common sizes and weights have
been fixed, the larger of the two thresholds works in every column.
\label{cor:frames-5}
\end{proof}

\begin{proposition}[The prime-sum gap at a fixed point]\label{prop:fixed-point-slope-gap}
Assume \eqref{eq:division-slopes}--\eqref{eq:upper-slopes} on one
family of columns, the reference bounds \(-\mu(M_u)\le A_-\),
\(\mu_{\max}(M_u)\le A_+\), and the level and exponent bounds
of Lemma~\ref{lem:exponent-level-spread}. Suppose the heights
\(h_u+S\log T_u\) have spread at most \(c_HB_0\), as in
Lemma~\ref{lem:division-height-spread}. Put \(X=x^{c_{HH}}\),
\(E_0=c_E m B_0\), and \(V=(c_H+2c_E m)B_0+S c_{\rm count}\log k\).
With the height cutoff \(H=bX\), the normalized Gram and slope
hypotheses of Lemma~\ref{lem:frame-slope} imply
\begin{align}
 S_{\mathrm{prime}}-S_{\mathrm{cent}}\le{}&
 A_-+E_D+E_0+SQ_c+\overline L+L_P+\log|W| \notag\\
 &+\frac{4\log(1+X)(A_++A_-+V+E_1+\overline L+b)
               +(\eta+4\sqrt e)Xb}{1-\delta}.
 \label{eq:pointwise-slope-comparison}
\end{align}
Here \(\overline L\) uses weights proportional to \(\rk N_u\).
\end{proposition}

\begin{proof}
\proofstep{The mean of the adjusted slopes.}
Put \(s=\overline{\mu(N_u)-h_u}\). Eliminating \(F\) between
\eqref{eq:division-slopes} and \eqref{eq:lower-slopes}, and using
the preceding identity, gives
\begin{equation}\label{eq:mean-slope-lower}
 s-\mu_C\ge G_a-S_{\mathrm{cent}}
           -(A_-+E_D+E_0+SQ_c+\overline L).
\end{equation}
\proofstep{The largest adjusted slope.}
To control the upper separation, subtract
\(\mu(N_{u'})-h_{u'}\) from \(\mu_{\max}(N_u)-h_u\).
The common terms \(G_a+F\) cancel in
\eqref{eq:lower-slopes} and \eqref{eq:upper-slopes}. Write
\(H_u=h_u+S\log T_u\) and use the height spread of
Section~\ref{sec:word-division}, say \(|H_u-H_{u'}|\le c_H B_0\).
Then
\[
 (-E_u-h_u)-(-E_{u'}-h_{u'})
  =(E_{u'}-E_u)+(H_{u'}-H_u)+S(\log T_u-\log T_{u'}).
\]
The three differences are bounded by \eqref{eq:slope-word-count-exponent}
and the height estimate, since \(S\ge0\). The required
spread is at most \(V=(c_H+2c_E m)B_0+S c_{\rm count}\log k\).
The difference of slopes is thus at most
\(A_++A_-+V+E_1+L_{u'}\). \label{cor:frames-6}
Average over \(u'\) and maximize over \(u\) to obtain
\[
 U=\max_u(\mu_{\max}(N_u)-h_u)-s
          \le A_++A_-+V+E_1+\overline L .
\]

\proofstep{Integrating the slope tail.}
Let \(X=x^{c_{HH}}>0\) and use \(H=bX\) by \eqref{eq:marked-frame-height}.
Concavity of \(\log(1+y)\) gives its tangent bound at \(X\).
Thus for \(U>0\)
\[
 U\log(1+H/U)
 \le U\log(1+X)+\frac{H-UX}{1+X}
 \le (U+b)\log(1+X).
\]
The inequality holds at \(U=0\) by continuity.
Apply the slope upper bound, allowing the additional
\(\log|W|\) term in the source, and combine it with
\eqref{eq:mean-slope-lower}. This gives \eqref{eq:pointwise-slope-comparison}.
\end{proof}

The construction must give the same column lattice data for
\eqref{eq:division-slopes} and \eqref{eq:lower-slopes}, with
uniform bounds for all terms in \eqref{eq:pointwise-slope-comparison}.

Section~\ref{subsec:fixed-weight-absorption} bounds the right side by the bracket in Lemma~\ref{lem:comparison}. \label{cor:frames-7}

\subsection{Small shares at the common cutoff}
\label{subsec:small-share-slope}

The pointwise comparison \eqref{eq:pointwise-slope-comparison} still contains an error from primes at which too small a proportion of the column survives. Estimate this contribution before summing the remaining prime losses. Write \(Q=|T|mb\), \(a=xb\), and fix a
point carrying the common weight vector. For the set
\(\mathcal S\) of candidate primes, let \(\sigma_p\ge0\)
be the share statistic and \(\theta\) its threshold. The small-share condition is
\begin{equation}\label{eq:gram-small-share-input}
 \sum_{p\in\mathcal S}\frac{\log p}{p}
        \bigl(2\mathbf1_{\sigma_p>\theta}+4\sigma_p\bigr)
       \le x^{-9}.
\end{equation}
\begin{lemma}[Three losses at the small-share cutoff]\label{lem:three-share-losses}
At one point assume \eqref{eq:gram-small-share-input}. Suppose the
local image loss satisfies \(i_p\le 4\sigma_p a m|T|/p\), the
second-class loss is bounded by the corresponding indicator term
with multiplier \(2am|T|\), and the loss for missing orbit labels
is at most \(4am\sum_{p\in\mathcal P_R}|B_p|(\log p)/p\).
Assume also the bad-prime count
\[
 \sum_{p\in\mathcal P_R}|B_p|\frac{\log p}{p}\le |T|x^{-2+\epsilon}.
\]
Then the image loss summed over the candidate set \(\mathcal S\),
the loss above the share threshold in the second class, and the loss
of labels without the required orbit condition are respectively at most
\[
 Qx^{-8},\qquad 2Qx^{-8},\qquad 4Qx^{-1+\epsilon}.
\]
\end{lemma}
\begin{proof}
Let \(i_p\) be the loss in the image of the local map. Its pointwise
bound \(i_p\le4\sigma_p a m|T|/p\) gives
\[
 \sum_{p\in\mathcal S}(\log p)i_p
             \le a m|T|x^{-9}=Qx^{-8}.
\]
For the primes in the second class, membership in
\(\mathcal S\) and \(\sigma_p>\theta\) give the loss
bound \(2a m|T|x^{-9}=2Qx^{-8}\) by the indicator
term in the same sum. Both estimates hold on each shift
for the same share statistic.

At a prime \(p\) with a selected orbit, every index with larger
weight that fails the local orbit condition belongs to the bad
set \(B_p\). The occurrence of any such index puts \(p\) in the
large-order residue set \(\mathcal P_R\).
The common bad-prime count is
\[
 \sum_{p\in\mathcal P_R}|B_p|\frac{\log p}{p}\le |T|x^{-2+\epsilon}.
\]
Indeed this is the count over the original set of indices with larger weights,
before passing to the subset of indices without the local orbit condition for a
particular record. The loss of those labels is bounded by
\(4am\) times that sum, which is \(4Qx^{-1+\epsilon}\).
\end{proof}
Section~\ref{subsec:loss-prime-account} derives the remaining
contributions from the orbit gains and the same size bounds for the auxiliary data.
\label{cor:frames-8}

\subsection{The prime sums at a selected point}\label{subsec:loss-prime-account}
\hypertarget{source-prime-class-loss}{}

For primes outside the small-share estimate the loss depends on the Frobenius orbits and their endpoint weights. Both sums below use the translate already fixed in the slope comparison.

\begin{established}[Cutoffs and widths in the prime-loss sum]\label{hyp:selected-prime-loss-weights}
Fix the target character family and \(0<\epsilon<1/52\).
Use a completed translate of Proposition~\ref{prop:translate-completion}, at
scale \(x\), with \(n=|A|\), real degree \(f=2n\) and last rank
\(N_m=k+m(t+1)\). Write \(a_s,-b_s\) for its two endpoint entries at
\(s\in T\), \(i_p\) for the local image loss, and \(G_a\) for the
arithmetic quantity in the local lattice comparison. Set
\[
 \begin{gathered}
 z=x^\epsilon,\quad b=e^z,\quad a=xb,\quad X'=e^{z^{1/3}},\quad X''=xnX',\\
 b'=a/X',\quad w''=a/X'',\quad \delta_1=z^{-1/8},\quad \ell=\log x .
 \end{gathered}
\]
Then the prime split \(P=\lceil k^2x^{26}\rceil\), the relative width
\(h=1/(kx^{13})\), and the share cutoff \(1/z\) may be used at every
translate with \(a\le a_s\le a(1+h)\) and \(b_s\ge b'\).
\end{established}
\begin{proof}
The independent integer choices in Section~\ref{subsec:frobenius-selection}
place the first entry in \([a,a+a/(kx^{13})]\) and the second in an
interval starting at \(b'\). Passing to the completed translate retains
these entries. Squaring the inverse first width gives precisely the
prime split \(P=\lceil h^{-2}\rceil\). All implied constants below are
bounded by a common \(K\ge1\) after the family and \(\epsilon\) are fixed.
\end{proof}

\begin{proposition}[Prime loss before and after the lattice map]
\label{prop:joint-prime-loss}
With the two weight functions, the prime cutoff \(P\) and the
share threshold \(z^{-1}\) of Lemma~\ref{hyp:selected-prime-loss-weights},
there is a constant \(K\) for the fixed
family and \(\epsilon\) such that eventually, uniformly in
the translate index,
\begin{equation}\label{eq:joint-prime-loss}
 S_{\mathrm{prime}}-G_a\le(3K+13)m|T|b,\qquad
 \sum_p(\log p)i_p\le m|T|b .
\end{equation}
\end{proposition}

\begin{proof}
\proofstep{Move both sides to rational primes with the same local degree weights.}
For a lift \(j\) and a top \(s\in T\), let \(W_{j,s,p,i}\) be the
filled weight at the partner reached after \(i\) Frobenius steps.
The second weight \(b_s\) enters this filling on the top positions.
At an inert prime (the signal or the residue class in the host)
the contribution before its logarithmic weight is
\[
 A_{j,s}(p)=\sum_{i>0,\ i\ {\rm odd}}
                 (a_s p^{-i}-W_{j,s,p,i})_+ .
\]
Every filling is at least \(w''\ge1\). This follows at a top from
\(b_s\ge b'\) and the positive offset. The lower offset is \(b\)
at the other positions with larger weights and \(w''\) at positions
with smaller weights. Thus at the positions with odd exponents one need retain only
\[
 1\le i\le\lfloor\log_p a_s\rfloor+1,\qquad i\equiv1\pmod 2 .
\]

At \(d\) there is one sum with all
positive exponents and the discriminant-partner weights. These
weights are nonnegative and beyond the indicated cutoff the partner lies in
the supported weights, where the weight is at least one.
The same truncation applies there. In particular the weight sum at the selected point has the cutoff \(\lceil2a\rceil\) of the selector.

Write \(g_j(p)\) for the gain kept at lift \(j\) and \(\tau(p)\) for the baseline cost of the restricted lattice; the total kept gain at the prime is
\[
 \sum_{j=1}^m g_j(p).
\] The point identifies
each local gain, cost and set of top indices on selected orbits with the corresponding quantity indexed by the rational prime. Inert primes at most \(X''\) have every place above them in
the support. For larger primes both gain and baseline cost vanish,
and at a non-inert place their difference contributes zero.
Indeed the candidate-orbit set is empty above \(X''\), and the
definition of the gain cuts it off there. If \(F\) is the
coefficient field then 
\[
 \sum_{w\mid p}e_w f_w=[F:\mathbb Q],\qquad \frac1{[F:\mathbb Q]}\sum_{w\mid p}e_w\log\mathrm Nw=\log p.
\]

Consequently the sum of the normalized weights
\(e_w\log{\rm N}w/[F:\mathbb Q]\) above \(p\) is \(\log p\).
Hence the ordinary part of the difference in
\eqref{eq:joint-prime-loss} is exactly
\begin{equation}\label{eq:loss-deficit-per-prime}
 2\sum_p \log p\,
 \left(\sum_{j,s} A_{j,s}(p)-\sum_j g_j(p)+\tau(p)\right),
\end{equation}
over inert primes. The \(d\)-sum has multiplier \(\log d\).
The prime assigned to the sequence cannot occur in this inert sum. Its chosen norm
generator and congruence give \(T_u\equiv1\) modulo the conductor of the auxiliary field times \((g_0g_1)^2\). In particular it is prime to
the conductor and \(T_u\equiv1\bmod 4d\), so quadratic reciprocity
gives \(({-d}/T_u)=1\). The dyadic prime divides the auxiliary conductor. Its splitting in the quadratic subfield follows from the discriminant test:
\[
 -d\equiv1\pmod 8\quad\Longrightarrow\quad 2\text{ splits in }\mathbb Q(\sqrt{-d}).
\] These observations
justify the inert drop rule without an exception depending on the sequence.
\label{cor:prime-class-loss-1}

\proofstep{Charge the top indices orbit by orbit.}
Let \(S_c\) be the weight bound in the auxiliary data, with
\(c_G=4S_c\), and let \(\eta\) bound the weight offset.
Write \(\lambda_{\rm ord}\) for the order parameter of
Section~\ref{sec:host-edge-tests} and \(r_G\) for the plane
dimension, fixed for this choice of \(\epsilon\).
Let \(C\ge1\) be a constant that may be increased a fixed
number of times before \(x\) is chosen. Define
\[
 \begin{gathered}
 K_0=k^2x^{78},\quad L_p=1+\lceil\log_p K_0\rceil,\qquad
 H_\circ=\exp(x^{\epsilon/2}),\qquad S_k=S_c+k+\eta,\\
 Q_1=C r_G b'\ell,\qquad
 Q_2=C\bigl(b'/x+\eta+S_c+(w''+t)\ell/\delta_1\bigr)+H_\circ,\\
 O_p=C\ell\left(\frac{\lambda_{\rm ord} a}{pk}
       +\begin{cases}a p^{-L_p}&p<P,\\ a/p^2&p\ge P,\end{cases}\right),\\
 B_p^*=C S_k\log_p(2a)+{\bf1}_{p<X'}
                \bigl(O_p+{\bf1}_{p<P}Q_1+{\bf1}_{p\ge P}Q_2\bigr),\\
 F_p=Cf\bigl(N_m^2\log_p a+(\log_p a)^2\bigr).
 \end{gathered}
\]
At a prime \(p\) retained in the sum, write \(T_p\) for the
top indices on its selected orbits. We claim
\begin{equation}\label{eq:top-service-prime}
 \sum_j\sum_{s\in T_p}\big(A_{j,s}(p)-B_p^*\big)-m F_p
                  \le \sum_jg_j(p)-\tau(p).
\end{equation}

To verify this, first suppose that the set \(V\) of selected orbits
at \(p\) is nonempty. It consists of the good orbits of a candidate
prime in the residue class below \(X'\). Number one of them cyclically
with length \(n'\). Write \(v\) for its number of tops and fix \(j\).
By the good-orbit count in Section~\ref{sec:host-edge-tests},
\(v\ge1,\ n'\delta_1\ll v\ell,\ n'\ll v\ell/\delta_1\), and
\(N_+=40n'\delta_1\). Write \(L=\log_p(2a)\). The local digit cost is
\[
 d_v=c_1 n' N_m^2L+c_2 vL^2.
\]

For each top set \(M_s=\sum_{i>0,\ i\ {\rm odd}}
\min(a_sp^{-i},W_{j,s,p,i})\).
Termwise addition gives \(A_{j,s}(p)+M_s=a_s p/(p^2-1)\).
With \(g^U_{jv}\) the gain using the up operator, this gives
\begin{equation}\label{eq:orbit-min-balance}
 \sum_{s}A_{j,s}(p)+\sum_{s\ {\rm even\ top}} M_s
                  \le g^U_{jv}+d_v+3v c_G L .
\end{equation}
In fact the first term of the up gain uses \(a_s-2c_G\) at all
tops.

At an odd top it also subtracts the minima with the adjusted
capacity \(W_{j,s,p,i}-\kappa_j+c_G\), where
\(0\le\kappa_j\le N_m\), and this capacity is nonnegative.
Writing the difference as a sum of positive parts, the two separate costs are bounded by
\[
 2c_G\sum_{i\ {\rm odd}}p^{-i}\quad\text{(numerator)},\qquad c_G\quad\text{per originally positive term (capacity)} .
\] There are at most \(L\) such indices since \(W\ge1\).
At an even top only the numerator shift is needed. Here \(L\ge1\).

Let \(P_{jv}\) be the exponent of \(p\) in the Hecke operator identity \eqref{eq:tag-p5}, \(t_v\) the real order
and \(D_p\) the integral depth in the local edge test. For \(p<P\)
the exponent \(P_{jv}\) is at most
\(3r_G[v b'+n'(w''+\eta+c_G+N_m)]\).
The expression for \(t_{\rm II}\) in Section~\ref{sec:host-edge-tests}
costs
\[
 O\!\left(r_G v
    (\lambda_{\rm ord} a/(pk)+a p^{-L_1}+c_0+k p^{G'})
                  +r_G n'c_G\right).
\]
The local exponent here is \(L_1=\lceil\log_p((kx^{26})^4)\rceil\).
Since \(p<k^2x^{26}\), \(pK_0\le(kx^{26})^4\), so \(L_p\le L_1\).
Also \(c_0=2b'+2c_G=O(b')\), \(k p^{G'}\le b'\) by the
small-prime bound there, and
\((\ell/\delta_1)(w''+\eta+c_G+N_m)\ll b'\).

For the last estimate \(w''=b'/(xn)\), the ranks are polynomial
in \(x\), and \(\log(1+\eta+c_G)=O(x^{\epsilon/4}+\ell)\).
In particular the displayed formula for \(D_{\rm II}\) gives
\[
 D_{\rm II}\ll r_G v b'\ell .
\] Indeed the factor multiplying
\(N_+\log_p a\) is polynomial, \(N_+\ll v\ell\), and
\(n'c_G\ll v b'\). Together these estimates give \(P_{jv}+\lceil t_v\rceil+D_p+1\le v(Q_1+O_p)\). Integer roundings cost at most
one apiece.

For \(P\le p<X'\) there is a term in \(t_{\rm I}\) of the size of
a top weight. Discarding this term as an error would exceed the
available bound. At each odd top \(y\) take the preceding even top
\(u\) and its positive distance \(i\), with summand zero if no
such top exists. The length \(n'\) is even on these candidate
orbits, so \(i\) is odd also across the end of the cycle.

The pair \((u,i)\) recovers the residue \(y\), proving injectivity:
\[
 (u,i)=(u',i')\quad\Longrightarrow\quad y\equiv u+i=u'+i'\equiv y'\pmod{n'} .
\]
Thus the sum of the terms \(\min(b_y,a_up^{-i})\) costs at most
\(\sum_{u\ {\rm even\ top}}M_u\), as \(W_y\ge b_y\). The variable
part of the exponent \(P_{jv}\) is
\(\sum_{y\ {\rm odd\ top}}(W_y-\kappa_j-b_y)\).
The chosen offsets bound this by \(v(b'/x+\eta)\).
This explains the surviving term \(b'/x\) in \(Q_2\).

The residual terms in the large-prime order are bounded, up to fixed
constants, by
\[
 vN_1\big(a/p^2+\lambda_{\rm ord} a/(pk)\big)
  +v\big(p^{G'}c_6+kp^{G'}+k\log_p a+c_G\big)+2n'c_G .
\]
Here \(c_6=8(D_{\rm I}+1)+O(N_+k\log_p a+(N_++n')c_G)\)
is the internal conductor. The remaining contribution to \(P_{jv}\) is
\(O(n'(w''+\eta+c_G+N_m))\).

We have \(N_1\ll\ell\), and the \(n'w''\) part costs
\(O(vw''\ell/\delta_1)\).
The depth \(D_{\rm I}\) is polynomial in \(x\). Moreover
\(\log(p^{G'})=O_\epsilon(x^{11\epsilon/24}\ell)\).
Set aside the two terms already bounded by \(O_p\) and \(w''\ell/\delta_1\). The remaining terms per top, including conductors, depth and roundings, have logarithm bounded by
\[
 O_\epsilon\!\left(1+x^{11\epsilon/24}\ell+x^{\epsilon/4}+\log x\right)=o(x^{\epsilon/2}).
\] Their sum is therefore at most \(H_\circ\).
The coefficient of \(H_\circ\) can be one for all sufficiently late
points. This proves from \eqref{eq:orbit-min-balance}, in both
ranges, that
\[
 \sum_{s\text{ on }v}(A_{j,s}(p)-B_p^*)\le
     g^U_{jv}-P_{jv}+d_v-(\lceil t_v\rceil+D_p+1).
\]
These are the same orders and depths that define the common baseline
of the restricted column in Section~\ref{subsec:column-orbit-length}.

Sum over \(v\in V\) and the \(m\) lifts. The kept gain uses
\(\max(g^U_{jv}-P_{jv},g^0_{jv},0)\) on \(V\), and
\(\max(g^0_{jv},0)\) elsewhere. The restricted-lattice baseline at \(p\) is
\[
 \tau(p)=e_0+m\sum_{v\in V}(\lceil t_v\rceil+D_p+1).
\] The Frobenius cosets partition the relevant labels, using both
signs at signal primes. Thus the sum of their lengths is at
most \(f\) and the number of tops in them is \(|T|\le f\).

In conjunction with \(p\le X''\le a\), \(L\le2\log_p a\)
and \(e_0\ll fk^2\log_p a\) from \eqref{eq:tag-c2}, the bounds defining
\(d_v\) give \(\sum_{v\in V}d_v+e_0\le F_p\).
The sum of the negative baseline terms is
\(-\tau(p)+e_0\). Consequently the orbit estimates give
\[
 \sum_j\sum_{s\in T_p}(A_{j,s}(p)-B_p^*)\le
  \sum_j g_j(p)-\tau(p)+m\sum_{v\in V}d_v+e_0
  \le \sum_j g_j(p)-\tau(p)+mF_p .
\]
The last step uses \(m\ge1\) to cover the single \(e_0\)
with the same allowance.

If \(V\) is empty, \(\tau(p)=0\). Below \(X'\) in the residue
class no top indices are then selected. On a signal prime or above the
cutoff the base gain suffices. In the base formula of
Section~\ref{subsec:hecke-entries} the excluded odd displacements
are first visits to a different top on the unbarred side.
At a signal prime the odd partner is on the barred side. For \(p\ge X'\), write \(U_i\) for an excluded numerator
and \(W_i\) for its fill at the destination. Then
\[
 \frac{U_i}{p^i}\le\frac{a(1+h)}{X'}
       \le b'+\frac{b'}{x}\le W_i,
 \qquad \bigl(U_i p^{-i}-W_i\bigr)_+=0 .
\]
The positive-part comparison used for \eqref{eq:orbit-min-balance}
bounds the shifts by \(3c_G\log_p(2a)\) per top, and the digit
costs sum to \(F_p\).

In the residue class with zero Frobenius step, above this cutoff
every partner is the starting top and every required positive
part is already zero. This covers the convention that the gain
in that case is zero. For primes near \(X''\) where
\(\log_p(2a)\ge1\) the same shift bound applies.
The calculation has used the unshifted fill of the actual prime sum.
Replacing the actual fill \(W\) in one summand by \(W-\kappa_j\)
has the elementary cost
\[
 0\le (v-W+\kappa_j)_+-(v-W)_+\le\kappa_j\le N_m .
\]
Here \(W-\kappa_j\ge1\), so for \(v=a_sp^{-i}\) the
difference can be nonzero only when \(p^i<a_s\). Summing over
the allowed \(i\) gives at most \(N_m\log_p(2a)\). Increase the
fixed constant \(C\) in the \(S_k\) term to absorb this cost.
\label{cor:prime-class-loss-2}

\proofstep{Sum the contributions from the selected orbits, exceptional primes, and the large-prime tail.}
Let \(\mathcal D\) be the remaining dropped inert primes, and
write \(\mathcal H\) for the candidates with share strictly above
\(z^{-1}\). The three sets used in the summation are disjoint:
\begin{equation}\label{eq:inert-dropped-prime-classes}
 \begin{aligned}
 \mathcal D
  ={}&\{p\in\mathcal D:p>X''\}
    \ \sqcup\ \{p\in\mathcal D:p\le X'',\ p\in\mathcal H\}\\
   &\ \sqcup\ \{p\in\mathcal D:p\le X'',\ p\notin\mathcal H\}.
 \end{aligned}
\end{equation}
On all three \(g_j(p)=\tau(p)=0\), and the set of selected orbits is empty.
The candidate test requires a residue-class prime below
\(X'\) of large order and excludes powers in the forbidden band.
The share test in \eqref{eq:inert-dropped-prime-classes} decides whether a candidate's good orbits
enter the selected set. Candidates with share strictly above
\(z^{-1}\) remain in the second dropped class and in the weighted
sum below. For an odd-order residue step the
projection to the dyadic quotient is zero, so it also fails the
large-order condition on these hosts. The drop and no-gain tests
agree on inert primes since every prime assigned to the sequence splits.

Let \(Q=m|T|b\). For every \(s,j\),
\[
 2(\log p)A_{j,s}(p)
     \le 2a_s(\log p)\frac{p}{p^2-1}\le4a\,\frac{\log p}{p}.
\]
Apply this estimate to \(T\setminus T_p\) at a prime with a selected orbit.
There is no missing top at a signal prime or at \(p\ge X'\).
Otherwise a missing top lies in a bad cycle: the prime is a
candidate in the large-order residue set, and its retained tops
are exactly those with good representatives. The weighted
bad-top count from Section~\ref{subsec:frobenius-selection} then
gives \(4Qz/x\) as in Section~\ref{subsec:small-share-slope}.
At dropped primes of large share the same inequality gives, after summation, the contribution bound
\[
 2Qx^{-8}.
\]

At any other dropped \(p\le X''\) the definition of the drop set
gives at least one of the following.
\begin{itemize}
\item The prime is at most \(M=\max(X'(1+C_r\rho),X'')\) and its
order in \((\mathbb Z/2^{e_2+2}\mathbb Z)^\times/\{\pm1\}\)
is at most \(J=\lfloor x^{\epsilon/4}\rfloor\). Here \(e_2\) is
the exponent of the dyadic modulus and \(C_r\) is fixed. The Frobenius
projection identifies this order with that of the dyadic component
of the Frobenius element, so its order divides that of the element. The dyadic component has order a power of two. For a Frobenius element of odd order it therefore satisfies
\[
 \operatorname{ord}(\text{dyadic component})\mid\gcd(2^a,\operatorname{ord}(\text{Frobenius element}))=1,
\]
with \(2^a\) the order of the dyadic quotient.
\item \(p\ge P\) and \(X'(1-\rho)\le p^i\le X'(1+\rho)\) for some
integer \(i\ge1\), where \(0\le\rho\le x^{-2}\) is the common
relative width from the host construction.
\end{itemize}
The alternatives follow by removing the high and large-share cases
from the no-gain test. Any remaining residue-class obstruction is
a short order below the enlarged cutoff, a forbidden power, or the
odd-step case. On the signal class those residue conditions are absent.

In the first alternative the dyadic calculation in Section~\ref{subsec:frobenius-selection} leaves only two classes:
\[
 p\equiv\pm1\pmod H,\qquad H=2^{e_2+2-\lfloor\log_2 J\rfloor}.
\]
Here \(x^2\le 2^{e_2}\le n\), so \(H\ge4\) eventually.
Summing \((\log p)/p\) over \(p=jH\pm1>1\) up to \(M\) gives at most
\[
 C_{\rm ord} J\,\frac{\log^2(2^{e_2}M)}{2^{e_2}}
\]
by the harmonic sum and \(H\ge 2^{e_2}/J\), for an absolute
\(C_{\rm ord}\).

In the second alternative \(p\ge8\) and
\(p^i\le2X'\), so \(i\le\lfloor\log X'\rfloor\) at large \(x\).
Writing \(R=(X')^{1/i}\) confines \(p\) to
\([R(1-\rho),R(1+\rho)]\). The count of integers is at most
\(2\rho R+1\) and \(p\ge\max(P,R/2)\). After extracting \(\log(2X')\), the interval length and one possible endpoint contribute respectively \(4\rho\) and \(1/P\). Together the resulting contribution is bounded by
\[
 \left(4\rho+\frac1P\right)\log(2X').
\]
Multiplying each of these two prime sums by \(4m|T|a\) bounds
the losses even when the alternatives overlap.

If \(p>X''\), \(A_{j,s}(p)\) can be positive only for \(i=1\) and
\(p<X''(1+h)\). Indeed, the scale and weight bounds together imply
\[
 \frac{a}{p^2}=\frac{w''X''}{p^2}\le \frac{w''}{2}\le\frac{W}{2}.
\] The remaining positive part is at most \(w''h\).
There are at most \(X''h+1\) such integers, giving a cost at most
\(2m|T|\log(2X'')(ah^2+w''h)\). At \(d\) the exponents include
the even ones, and summing the full geometric series gives
\(a_s/(d-1)\le 2a/d\) for \(d\ge4\). The cost is at most
\(2m|T|a\log d/d\). Eventually \(d\ge x^2>e\). Since \(t\mapsto(\log t)/t\) decreases for \(t>e\),
\[
 \frac{\log d}{d}\le\frac{\log(x^2)}{x^2}=\frac{2\ell}{x^2}.
\]

It remains to sum \(F_p,B_p^*\) from \eqref{eq:top-service-prime}.
There are at most \(X''\) primes with selected orbits. Use
\(\log p\log_p a=\log a\) and \(1/\log p\le2\). This gives the
following two bounds for the floor and common per-top costs,
respectively,
\[
 2m C fX''\big(N_m^2\log a+2(\log a)^2\big),
 \qquad 2m|T|C S_kX''\log(2a).
\]
For \(P\) and \(X'\) the elementary Chebyshev inequality
\(\sum_{p\le y}\log p\le y\log4\) gives a combined bound
\(2m|T|\log4\,(P Q_1+X'Q_2)\).

For the order term use
\[
 p^{1-L_p}\le K_0^{-1},\qquad
 \sum_{p\le y}\frac{\log p}{p}\le\log y+\log4 .
\]
For the second inequality, compare
\(\sum_{p\le n}\lfloor n/p\rfloor\log p\le \log(n!)\)
with \(n\log n\) and apply Chebyshev to the rounding error,
with \(n=\lfloor y\rfloor\). Also
\(\sum_{j\ge P}(\log j)/j^2\le
2(\log(P-1)+1)/(P-1)\) by integral comparison for \(P\ge3\).
The order term therefore costs at most
\[
 2m|T|C\ell a\left[
  (\lambda_{\rm ord}/k+1/K_0)(\log X'+\log4)
       +2\frac{\log(P-1)+1}{P-1}\right].
\]

\proofstep{Separate the one term that does not vanish after division by \(Q\).}
Use \(K\ge\max(1,C,C_{\rm ord})\) fixed and put
\(s_x=100K z^4/x\). Eventually \(s_x\le1\).
All rank and cutoff bounds in these comparisons are simultaneous:
\(x^2\le 2^{e_2}\le n,\ x^2\le k,\ x^4+1\le P\);
\(n,k,N_m,t,P\le x^{E}\) with \(E\) fixed for this family;
\(\lambda_{\rm ord}\le z^2\), \(S_c,\eta\le \exp(z^{1/4})\),
and \(1\le X'(1+C_r\rho)\le 2X'\).

The two factors in \(S_c=4k\Delta_{\rm host}\) each contribute at most half the exponent on the eventual scale, since \(k\) is polynomially bounded:
\[
 \log S_c=\log(4k)+\log\Delta_{\rm host}\le z^{1/4}/2+z^{1/4}/2=z^{1/4}.
\] The stated lower
bound on \(d\) follows from the positive power lower bound on
\(\log d\) for this fixed family. One may take
\(E=\max(40,C_t+3)\) using the size choices in
Section~\ref{sec:family-window}. These bounds apply before any
choice of \(j,s,p\) or translate.

Divide the preceding costs by \(Q\). The two order and band
exceptional sets cost respectively
\(O(K z^{1/4}(\ell+\log X')^2/x)\) and
\(O((\log X')\log(2X')/x)\) since \(n\le x^E\).
The high-prime and discriminant bounds are
\(O(\log(2X'')/x)\) and \(O(\ell/x)\) since \(h\le1/x\) and
\(w''\le b\).

Since \(|T|\ge1\), the floor and common per-top costs are bounded by the following quantity, whose exponential decay comes from \(b=e^z\):
\[
 O\!\left(\frac{x^{O(1)}X'\exp(z^{1/4})}{b}\right)=\exp\bigl(-z+z^{1/3}+z^{1/4}+O(\log x)\bigr).
\]
The small-prime contribution \(P Q_1/b\) is \(O(x^{O(1)}/X')\).
Each expression is \(o(s_x)\). For the order term for the selected orbits,
\(\lambda_{\rm ord}/k+1/K_0\le(z^2+1)/x^2\),
\(\log X'=z^{1/3}\) and \(\log(P-1)=O(\ell)\). Thus its
ratio to \(Q\) is \(O(C\ell(1+z^2)(1+z^{1/3})/x)=o(s_x)\).

There is one main term. The contribution from the larger primes is
\[
 \frac{X'Q_2}{b}= C+
 C\frac{X'(\eta+S_c)}{b}
 +C\frac{\ell}{n\delta_1}
 +C\frac{X't\ell}{b\delta_1}
 +\frac{X'H_\circ}{b}.
\]
Here \(X'b'/x=b\) and \(X'w''=b/n\). Every term after \(C\)
is \(o(s_x)\), including the last one since
\(\log(X'H_\circ)=z^{1/3}+z^{1/2}=o(z)\).
This calculation needs the cancellation of the preceding-even-top
minima before estimating the contribution from the larger primes.

For sufficiently large \(x\), each of the two floor and common costs,
small-prime contribution, order contribution for the selected orbits, two first-class costs, high-prime
tail and discriminant cost is at most \(s_xQ\). The error from top indices outside the selected orbits
and large-share costs are each at most \(s_xQ\). The contribution from the larger primes is at most \((2C\log4+3s_x)Q\).
Summing the contributions gives
\[
 S_{\mathrm{prime}}-G_a\le(2C\log4+13s_x)Q.
\]
The share error is absorbed since
\[
 Qx^{-8}\le s_xQ .
\] The weighted sum of the output over the coefficient
places reduces to the same sum over candidate primes by the local
degree formula already used in \eqref{eq:loss-deficit-per-prime}.
The joint local hypotheses and the small-share weight condition
carry both inequalities to each selected point. The common threshold
is a finite maximum of the bounds uniform in the point.
\label{cor:prime-class-loss-3}
\end{proof}


\begin{lemma}[One small-share vector at each scale]\label{lem:one-share-vector}
Fix \(0<\epsilon<1/52\). At each sufficiently large scale \(x\) in a selected window there is an admissible vector \(o\) for which the shares satisfy
\[
 \sum_{p\in\mathcal S}\left(2{\bf1}_{\sigma_p(o)>x^{-\epsilon}}+4\sigma_p(o)\right)\frac{\log p}{p}\le x^{-9}.
\]
This one choice of \(o\) works at every point of the scale carrying it.
\end{lemma}

\begin{proof}
To select the weight vector, take one reference point at \(x\)
and let \(\mathcal W\) be its nonempty finite set of admissible
vectors. At a candidate prime sum the local exceptional fractions
over its good Frobenius cycles, with the factor \(m\) for the lifts.
The two local estimates in Propositions~\ref{prop:host-small-prime-test}
and~\ref{prop:host-large-prime-test} give the mean bound of
Paragraph~\ref{par:top-weight-selection}, namely
\[
     \mathbb E_{o\in\mathcal W}\sigma_p(o)\le m|T|x^{-12}
                         \qquad(p\in\mathcal S).
\] Here the expectation is the uniform
mean and the candidate set is fixed before the vector is chosen.
For \(\theta=x^{-\epsilon}\) and \(s\ge0\),
\(2\mathbf1_{s>\theta}+4s\le(2/\theta+4)s\).

Since the candidates are primes below \(X'=\exp(x^{\epsilon/3})\),
the elementary estimate \(\sum_{p\le y}(\log p)/p=\log y+O(1)\) gives
\(\sum_{p\in\mathcal S}\log p/p\le2x^{\epsilon/3}\).
After summing the mean-share estimate, the mean of the required
statistic is at most
\[
 (2x^\epsilon+4)m|T|x^{-12}\cdot 2x^{\epsilon/3}.
\]

The size bounds are \(m\le x^{(1-\epsilon/8)/2+\epsilon/100}\)
and \(|T|\le2x^{2+8\epsilon+\epsilon/100-\epsilon/8}\).
The last display is bounded by \(24x^{-19/2+c\epsilon}\) with
\(c=9+1/3+2/100-3/16\). For \(0<\epsilon<1/52\) one has
\(-19/2+c\epsilon+1/4<-9\). Taking \(x^{1/4}\ge24\)
gives a mean at most \(x^{-9}\), so some vector attains the
small-share inequality at the reference point. The share statistic
and candidate set are invariant among points of that scale carrying
the chosen vector, by the Galois and scale identifications.
\label{cor:frames-9}

\end{proof}

\begin{proposition}[Simultaneous estimates on the translates]\label{prop:estimates-each-translate}
At each sufficiently large scale in a selected window one can use
the same weight pair for all translates. For every translate
\(u_0\) there is a completed point \(i(u_0)\) where the count,
slope, frame, division and loss estimates hold together.
At every such \(i(u_0)\) the admissible cutoff and the exponents in the division and loss bounds are
\[
 N=\lceil2a\rceil,\qquad(e_D,e_L)=(2,0).
\]

\end{proposition}
\begin{proof}
By Lemma~\ref{lem:one-share-vector} choose an admissible weight pair
at one reference point. Proposition~\ref{prop:translate-completion} realizes
this pair at every translate of the common field and binary sequence
assignment. On identifying the indices, the larger weights, the
dimensions and the share statistic at \(i(u_0)\) have their chosen
values. In particular the two terms at each prime give precisely the
small-share bound used in Lemma~\ref{lem:three-share-losses}.
Every larger weight is at most \(2a\) and every denominator weight
is at least one. Thus the choice of \(N\) applies to the same
\(S_{\mathrm{prime}}\) in the local loss and in the slope inequality.

Fix the original character weights and then choose the family in the
order of Lemma~\ref{lem:window-parameter-choice}. Take the maximum of the
finitely many uniform thresholds for its realization and the stated
estimates. At every later \(i(u_0)\),
Proposition~\ref{prop:sequence-division} and the joint loss bounds of
Proposition~\ref{prop:joint-prime-loss} now give the claimed exponents.
\label{cor:frames-10}
\end{proof}

\subsection{Combining the estimates with the fixed weights}
\label{subsec:fixed-weight-absorption}

The prime losses can now be inserted in the slope inequality. Keeping the weight vector fixed in this substitution is what permits the later average over translates.
\begin{established}[Numerical bounds on a completed member]\label{hyp:absorption-inputs}\label{cor:frames-11}
Write \(n_T=|T|\), \(n_2=|T_2|\), and \(A=n_Ta\).
Here \(f=2n\), \(B_0=fmb\,x^\epsilon\), and
\(N_c=n_2 k b x^{2\epsilon}/n^2\) is the rank contribution
from the comparison in Section~\ref{subsec:division-length}.
Define the factors \(y,y_D,y_L\), the expression \(\mathcal B\),
and its weighted version \(P\) in that order.
\[
 y=x^\epsilon,\qquad y_D=x^{e_D\epsilon},\qquad
 y_L=x^{e_L\epsilon},\qquad
 \mathcal B=A(1+m^2/k)+m(B_0+N_c),\qquad
 P=y_Dy_Ly\,\mathcal B .
\]
These factors are all at least one on the size range, which has
\(1\le b\le a\), \(m\le k\), \(mb\le a\), \(kb'\le a\),
\(X''\le a\), and \(L,L_a\le2y\).
There is a constant \(K\ge6\), independent of the member and translate for
this family, for which the following estimates hold beyond one scale.

In the notation of Proposition~\ref{prop:fixed-point-slope-gap} these estimates give
\begin{align*}
 S Q_c&\le S K L_a(1+m^2/k),&
 E_D&\le K\big(A L_a+m(B_0+N_c)y_D\big),\\
 E_0&\le K mB_0,& E_1&\le K mB_0,\\
 V&\le K\big(2mB_0+A L_a+mB_0+S\log k\big),&
 L_P&\le K n_Tmb\,y_L,\\
 \overline L&\le K^2 n_Tmb\,y_L
      + K^2 n_Tm^2(a/k+b')L+KX''+1 .
\end{align*}
With \(X=x^{c_{HH}}\) and the good-label dimension deficit
\(\delta=1-D_{\rm good}/D\), the auxiliary bounds are
\[
 S\le2A,\qquad \log|W|\le a,\qquad
 (\eta+4\sqrt e)Xb\le K y A,\qquad (1-\delta)^{-1}\le2.
\]
\end{established}
\begin{proof}
The count in Lemma~\ref{lem:exponent-level-spread} gives \(Q_c=O(L)\).
Its exponent estimate gives \(E_0\le c_E m B_0\). The maximal-slope
bound \eqref{eq:tag-c8} and the rounding estimate in
Proposition~\ref{prop:shifted-column-slopes} give the stated bound for
\(E_1\) on the same shifted lattice, as in the proof of
Proposition~\ref{prop:marked-frame-height}.

The estimate for \(E_D\) is Proposition~\ref{prop:sequence-division}, in the form \eqref{eq:selected-word-division}. The bound for \(V\) adds the
exponent and count differences in \eqref{eq:slope-word-count-exponent}
and the variation of the adjusted height in
Section~\ref{sec:word-division}. For \(L_P\), use the first
bound in Proposition~\ref{prop:joint-prime-loss}, equation~\eqref{eq:joint-prime-loss}. For \(\overline L\),
combine the weighted output bound in the same equation with
the contributions at the assigned prime in
Section~\ref{subsec:word-period} and the rounding error.
For the remaining slope terms take the two bounds for \(A_-,A_+\)
in Lemma~\ref{lem:reference-column-slopes}. Each top weight is at most
\(2a\) by Lemma~\ref{hyp:translate-weights}. The binary injection of
Lemma~\ref{lem:window-parameter-choice} and the selected rank range give
\[
 S=\sum_{t\in T}a_t\le2n_Ta,\qquad
 \log|W|\le(k-2)\log2\le k\le a
\]
eventually, since \(k\le x^5\) and \(b\le a\) grows faster than
every fixed power of \(x\). Finally, Proposition~\ref{prop:selected-rank-comparison}
and Lemma~\ref{lem:trace-at-height} give
\[
 \delta\le e^{-x^{e_z/7}}\le\tfrac12,\qquad
 (\eta+4\sqrt e)Xb\le6b\le6yA .
\]
Take \(K\) to be the maximum of six and the finitely many family constants.
\end{proof}

\begin{lemma}[Absorption at one completed point]\label{lem:absorption-at-point}
At a completed member use the estimates of Lemma~\ref{hyp:absorption-inputs},
with its constant \(K\), in the slope inequality
\eqref{eq:pointwise-slope-comparison}. Let \(g,h\ge1\) satisfy \(\log k\le gy\) and
\(\log(1+X)\le hy\). If \(y\ge200(1+K)^2hg\), then
\[
 S_{\mathrm{prime}}-S_{\mathrm{cent}}
      \le x^{(e_D+e_L+3)\epsilon}\mathcal B .
\]
The bound applies at every translate once the common thresholds
for the inputs have been passed.
\end{lemma}

\begin{proof}
\proofstep{Bounding the two parts of the slope inequality.}
Take \(g,h\ge1\) such that \(\log k\le gy\) and
\(\log(1+X)\le hy\). They are fixed by the power bounds
on \(k\) and \(X\) once \(\log x\le y\). In estimating the
terms of \eqref{eq:pointwise-slope-comparison}, \(P\)
dominates each of \(A\), \(mB_0\), \(b\), \(X''\), and
\(A(1+m^2/k)\). The quadratic term in Lemma~\ref{hyp:absorption-inputs}
has the sharper bound
\[
 n_Tm^2(a/k+b')\le2Am^2/k,\qquad
 \overline L\le(5K^2+K+1)P .
\] The sum of the terms on
the first line of \eqref{eq:pointwise-slope-comparison}
is therefore at most \((5K^2+14K+2)P\), by using
\(4KP,3KP,KP,4KP\) for the count, division, centered
exponent and lower reference costs, and \(KP\) for \(L_P\).

In the second line the sum multiplying \(4\log(1+X)\)
is at most \((5K^2+13K+2Kg+2)P\). In fact the two
reference costs total at most \(6KP\), the spread costs
\((5+2g)KP\), and \(E_1\) adds \(KP\).
Combining these bounds and the trace-product term gives
\[
 S_{\mathrm{prime}}-S_{\mathrm{cent}}
 \le \big(5K^2+16K+2+8h(5K^2+13K+2Kg+2)\big)yP
 \le 200(1+K)^2 hg\,yP.
\]
The last inequality follows by \(h,g\ge1\) and comparison
of the three coefficients of the polynomial in \(K\).
The threshold \(y\ge200(1+K)^2 hg\) now gives
\[
 S_{\mathrm{prime}}-S_{\mathrm{cent}}
 \le x^{(e_D+e_L+3)\epsilon}\mathcal B.
\]
This derivation is uniform in the sequence and translate.
With the stated values of \(B_0\) and \(N_c\),
\(\mathcal B\) is the bracket of
Lemma~\ref{lem:comparison} with \(a_E=a\).
On the common inputs \(e_D=2\) and \(e_L=0\), so this
gives \(C=5\) at every point where the stated estimates
hold. \label{cor:frames-12}

\end{proof}

\begin{corollary}[The translate bound on one completed family]\label{cor:raw-gap-completed-scale}
At a sufficiently late scale of the family in Proposition~\ref{prop:estimates-each-translate},
fix its weight vector and take the selected member over each \(u\in A_0\).
Use its base and refined rows as identified in
Proposition~\ref{prop:linked-comparison-rows}. With \(a_E=a\), the centered local errors of Definition~\ref{def:translate-errors} satisfy
\[
 \max_{u\in A_0}\sum_{j=1}^m\sum_{t\in T}\Delta_{j,t}(u)
 \le x^{5\epsilon}\!
 \left[n_Ta_E\left(1+\frac{m^2}{k}\right)
      +m\left(2nmbx^\epsilon+\frac{n_2kb x^{2\epsilon}}{n^2}\right)\right].
\]
\end{corollary}

\begin{proof}
For each \(u\), Proposition~\ref{prop:translate-completion} gives the
completion at the prescribed levels. The strict cohomology bound on its
original host passes to its good characters by Lemma~\ref{lem:strict-budget-after-cut}.
Proposition~\ref{prop:uniform-parent-remainders} gives the geometric bounds
on the two covers of its fixed core. Thus Lemma~\ref{lem:joint-family} and Proposition~\ref{prop:parent-dual-cost} bound respectively the paired transport and the signed index of the localized global dual at the \emph{same} point. Proposition~\ref{prop:sequence-division} inserts them in the division bound there. With the other numerical inputs, Lemma~\ref{lem:absorption-at-point} now gives
\[
 \sum_{j,t}\Delta_{j,t}(u)
 =S_{\rm prime}(u)-S_{\rm cent}(u)
 \le x^{(2+0+3)\epsilon}\bigl[A(1+m^2/k)+m(B_0+N_c)\bigr].
\]
The right side is independent of \(u\). Substitute \(A=n_Ta_E\), \(B_0=2nmbx^\epsilon\), and \(N_c=n_2kbx^{2\epsilon}/n^2\), then take the maximum over \(A_0\).
\end{proof}

\begin{auxlemma}[A single numerical threshold on a family of marked rows]
\label{lem:uniform-numerical-application}
Fix the outer parameters and \(\epsilon>0\). Choose the middle height
exponent \(C_H=c_{HH}\ge1\) of Proposition~\ref{prop:marked-frame-height}
before choosing the sequences. Let \(S\subset[1,\infty)\) be a set of scales.
Require one threshold \(X_{\rm in}\) with
\[
 \forall x\in S\cap(X_{\rm in},\infty),\quad
 \forall u_0\in A_0(x),\quad
 \exists\ \text{a marked row over }(x,u_0)
\]
satisfying the following conditions. Here \(A_0(x)\) is the kernel in
Definition~\ref{def:translate-scale} for the host chosen at \(x\).
The two signed runs have the common form coordinates of
Lemma~\ref{lem:selected-form-level-tie}, the linked comparison of
Proposition~\ref{prop:linked-comparison-rows}, and the weighted numerator
estimate of Proposition~\ref{prop:weighted-dual-cost} with its original
place assignment. The period, measure and size data on those rows satisfy
the hypotheses of Propositions~\ref{prop:selected-rank-comparison}, \ref{prop:marked-frame-height} and~\ref{prop:fixed-point-slope-gap}. The remaining input estimates in
Lemma~\ref{hyp:absorption-inputs} hold with a common constant.
Assume the selected size range, in particular \(k\le x^5\).

There exists \(X_{\rm num}\), depending only on these uniform data, such
that for every \(x\in S\) above \(X_{\rm num}\),
\[
   \max_{u_0\in A_0}\sum_{j,t}\Delta_{j,t}(u_0)
     \le x^{5\epsilon}
       \left[n_Ta\left(1+\frac{m^2}{k}\right)
             +m\left(2nmbx^\epsilon+
                     \frac{n_2kb x^{2\epsilon}}{n^2}\right)\right].
\]
The conclusion holds on all scales covered by the hypotheses,
including the endpoints of any closed interval contained in
\(S\cap(X_{\rm num},\infty)\).
\end{auxlemma}

\begin{proof}
\proofstep{Use the actual cut and rank.}
Fix \(x\in S\) above \(X_{\rm in}\), take any \(u_0\in A_0(x)\),
and choose its marked row from the hypothesis. Lemma~\ref{lem:parent-slice-lattices}
identifies its branch and product lattices. Proposition~\ref{prop:uniform-parent-remainders}
therefore bounds every tested subspace in those same coordinates.
Use the same row to evaluate the numerator. At the rescaled argument
\(x'=x/\kappa_{\rm sc}\) the integers required by the local twists
have the range
\[
     1\le a'<t\le\lfloor x^{400}\rfloor
                 \le\lfloor(x')^{401}\rfloor
           \qquad(x\ge\max\{2^{100},\kappa_{\rm sc}^{401}\}).
\]
The bound on \(t\) is the size calculation in
Section~\ref{sec:family-window}. The last inequality follows from
\(x^{400}\le(x/\kappa_{\rm sc})^{401}\). Thus each signed numerator is
charged within the original list at \(x'\). The indices in
Lemma~\ref{lem:joint-family} and Proposition~\ref{prop:parent-dual-cost}
are now those of this one member. With the prime losses of
Proposition~\ref{prop:joint-prime-loss}, Proposition~\ref{prop:sequence-division}
gives \((e_D,e_L)=(2,0)\).

\proofstep{Use the same height in the trace and dimension bounds.}
Let \(K\ge6\) bound the finitely many uniform constants in
Lemma~\ref{hyp:absorption-inputs}, after making this substitution.
Put \(g=5\), \(h=C_H+1\), and set
\begin{align*}
 B&=1+\log18+16C_H/\epsilon,\\
 X_1&=\max\{X_{\rm in},2^{100},\kappa_{\rm sc}^{401},
             20^{1/C_H},(\max\{1,2/\epsilon\})^{2/\epsilon}\},\\
 X_{\rm num}&=\max\{X_1,B^{16/\epsilon},
                        [200(1+K)^2gh]^{1/\epsilon}\}.
\end{align*}
Here \(X_{\rm in}\) includes the uniform thresholds of the cited
estimates once the outer and middle parameters have been chosen.
For \(v=x^{\epsilon/16}\ge B\),
\[
 x^{\epsilon/8}=v^2\ge(16 C_H/\epsilon)v+\log18
       \ge C_H\log x+\log18 .
\]
Write \(E=\exp(-x^{\epsilon/8})\) and \(\sigma=x^{-C_H}\). Hence
\(E\le\sigma/18\) and \(\sigma\le1/20\).
Proposition~\ref{prop:selected-rank-comparison} gives
\[
 \max\{\eta_0,\eta_1,\sqrt{e_0},\sqrt d\}\le E,\qquad
             D_*/\sum_u d_u\le\sigma .
\]
Lemma~\ref{lem:trace-at-height} therefore gives
\((\eta+4\sqrt e)H\le6b\) throughout \(0\le H\le bx^{C_H}\).
The exponent used in the dimension estimate is the exponent used
to truncate this very slope integral.

\proofstep{Absorb the constants uniformly.}
Finally \(y=x^\epsilon\ge\log x\), since
\(\log x=(2/\epsilon)\log(x^{\epsilon/2})\le
(2/\epsilon)x^{\epsilon/2}\le x^\epsilon\).
Thus \(\log k\le gy\) and \(\log(1+x^{C_H})\le hy\). Our choice
ensures \(y\ge200(1+K)^2gh\). Lemma~\ref{lem:absorption-at-point}
bounds the difference on the chosen row by
\(x^{(2+0+3)\epsilon}\mathcal B\). Its right hand side depends
on the base data at \(x\), fixed before the choice of \(u_0\).
Substitution of \(\mathcal B\) and maximization over \(u_0\) prove
the assertion.
\end{proof}

At each scale of Proposition~\ref{prop:all-window-applications} the marked
family supplies a member over every \(u_0\in A_0(x)\). Corollary~\ref{cor:raw-gap-completed-scale}
identifies the two comparison rows on each such member. Thus the maximum in
Lemma~\ref{lem:uniform-numerical-application} ranges over the entire kernel,
as required for \(E\) in Proposition~\ref{prop:translate-fourier}.
Lemma~\ref{hyp:translate-weights} supplies its two weight ranges. Translation averaging has mass one. The weight replacement costs \(2b/x\) for each of the \(n_Tm\) pairs, hence \(2Q/x\) in \eqref{eq:fourier-compare}. The two rank inequalities in Lemma~\ref{lem:comparison} follow on the selected size range from
\[
 \frac{k}{n^2}\le\frac{2x^{9/2+10\epsilon}}{x^{4+16\epsilon}}
     \le x^{1/2-\epsilon}\quad(x^{5\epsilon}\ge2),\qquad
 xm\le x^{3/2+\epsilon/2}\le x^{9/2+10\epsilon}\le k .
\]
With the stated density, \(m\)-range and \(\Sigma(x)\) bounds this gives \(G_{\mathcal K}(x)\le x^{1/2+9\epsilon}\) at that scale.

The common threshold just obtained can be used on an entire
interval. In Proposition~\ref{prop:all-window-applications} the
family is chosen before the threshold. Lemma~\ref{lem:disjoint-common-pair-windows}
then permits its left endpoint to exceed \(X_{\rm num}\) and
any prescribed bound. This is the passage from one completed
scale to \eqref{eq:uniform-raw-gap} in
Theorem~\ref{hyp:uniform-arithmetic-data}.

For the weaker square-argument criterion in Lemma~\ref{lem:square-window-transfer},
the proof of the corollary is used only at the translate selected in
Lemma~\ref{lem:maximizing-translate}. Both indices must still come from
the same marked member there.
